\documentclass{subfiles}
\begin{document}
    We now turn our attention to the following problem:
        given a set \(M\) of items, how can we efficiently store the items in \(M\),
        while also being able to immediately tell if an item \(i \in M\)?
        That is, is it possible to determine if \(i \notin M\) by simply looking at the data structure?

    By now, it should be clear that such a problem is easily solved through hashing,
        with the only issue concerning space. In fact, if the space available for the hast table
        is much less than that actually needed, the usual hashing approach is infeasable.

    Here, a slightly different approach is discussed. In more details,
        we will overview the solution proposed in \cite{Bloom1970} by \citeauthor{Bloom1970},
        and refer the interested reader to the paper for major details.
        We focus on what in \cite{Bloom1970} is referred to as ``Method 2'';
        commonly known as Bloom's filter in the literature.

    The idea is the following: we represent our hash table as a collection of \(N\) individually
        addressable bits, initially set to zero. If \(M\) is the set of items to store,
        then for each item \(k \in M\) we identify a group of bits \(b_{1}, b_{2}, \ldots, b_{d}\),
        by means of some hash procedure, and set these to one.
        Similarly, if \(i\) is some item we wish to check belongs to the set,
        then we compute the hash for \(i\) and check if all the bits (in the hash) are set to one.

    \begin{remark*}
        It should be clear that, if for a given item \(k\), 
            all the bits in its hashing are set to one, then there's no guarantee that \(k\) belongs to the set.
            On the contrary, if any of the bits in the hashing is set to zero,
            then we are sure \(k\) doesn't belong to the set.

        Stated differently, we allow some items to be falsely recognized as belonging to the set,
            to immediately discard those items that do not.
            We refer to these falsely recognized items as the \emph{allowable fraction of error}.
    \end{remark*}

    We now proceed with the analysis of Bloom's filter.
        Let \(\phi\) be the expected proportion of bits still set to zero after \(n\) items 
        have been stored, and let \(d\) be the fixed number of bits set for each.
        Then 
        \[
            \phi = (1 - \sfrac{d}{N})^{n}
        \]
        It follows that an item is falsely recognized if all \(d\) bits are set to one,
        which happens with probaility \(P = (1 - \phi)^{d}\).

        If we assume \(d \ll N\), and we take the log base 2 of \(\phi\) we get approximately 
        \[\begin{aligned}
            \log_{2} \phi & =  \ln (1 - \sfrac{d}{N})^{n} \log_{2} (e) \\ 
                    & = -n (\sfrac{d}{N}) \log_{2} (e).
        \end{aligned}\]

    Therefore
    \[
        N = -n\log_{2} (P) \frac{\log_{2} (e)}{\log_{2} (\phi) \log_{2} (1 - \phi) }.
    \]
\end{document}
